% Part: computability
% Chapter: recursive-functions
% Section: primes

\documentclass[../../../include/open-logic-section]{subfiles}

\begin{document}

\olfileid{cmp}{rec}{pri}
\olsection{अविभाज्य संख्या}

परिबद्ध संख्यकीकरण आणि परिबद्ध लघुतमीकरण यांमुळे नेहमी आढळणारी अनेक
फलने व संबंध आदिम पुनरावर्ती आहेत हे दाखवण्यासाठी आपल्याला बरीच साधने
मिळतात. उदाहरणार्थ, ``$x$ ने $y$ ला निःशेष भाग जातो'' हा संबंध विचारात
घ्या; तो $x \mid y$ असा लिहितात. $y$ ला~$x$ ने बाकी न उरता भागता येत
असेल, म्हणजे $y$ ही~$x$ ची पूर्णांक पटीत संख्या असेल, तेव्हा $x \mid y$
हा संबंध सत्य असतो. (ते घडत नसेल, तर आपण $x \nmid y$ असे लिहितो;
भाजक धन असेल तेव्हा याचा अर्थ $y$ ला $x$ ने भागल्यावर बाकी $> 0$ उरते.)
दुसऱ्या शब्दांत, एखाद्या~$z$ साठी $x \cdot z = y$ असेल तेव्हा आणि केवळ
तेव्हाच $x \mid y$. असा~$z$ अस्तित्वात असेल, तर $\leq y$ असलेला एक
साक्षीदार निवडता येतो. म्हणून एखाद्या $z \le y$ साठी $x \cdot z = y$
असेल तेव्हा आणि केवळ तेव्हाच $x \mid y$. त्यामुळे $=$ आणि गुणाकार
यांच्यापासून परिबद्ध अस्तित्ववाची संख्यकीकरण वापरून $x \mid y$ या संबंधाची
व्याख्या पुढीलप्रमाणे करता येते:
\[
x \mid y \defiff \bexists{z \leq y}{(x \cdot z) = y}.
\]
अशा रीतीने $x \mid y$ हा संबंध आदिम पुनरावर्ती आहे हे आपण दाखवले.

नैसर्गिक संख्या~$x$ ही $0$ किंवा $1$ नसल्यास आणि तिला फक्त $1$ ने व
स्वतःनेच भाग जात असल्यास तिला \emph{अविभाज्य} म्हणतात. दुसऱ्या शब्दांत,
अविभाज्य संख्यांचे वैशिष्ट्य असे की, जेव्हा $y \mid x$, तेव्हा एकतर
$y = 1$ किंवा~$y=x$. $x$ अविभाज्य आहे की नाही हे तपासण्यासाठी आपल्याला
फक्त सर्व $y \le x$ साठी $y \mid x$ आहे की नाही हे तपासावे लागते, कारण
$y > x$ असल्यास आपोआप~$y \nmid x$. म्हणून, $x$ अविभाज्य असेल तेव्हा आणि
केवळ तेव्हाच सत्य असणाऱ्या $\fn{Prime}(x)$ या संबंधाची व्याख्या
\[
\fn{Prime}(x) \defiff x \geq 2 \land \bforall{y \leq x}{(y \mid x \lif y
  = 1 \lor y = x)}
\]
अशी करता येते आणि म्हणून हा संबंध आदिम पुनरावर्ती आहे.

अविभाज्य संख्या $2$, $3$, $5$, $7$, $11$, इत्यादी आहेत. या अनुक्रमात
स्थानांक~$x$ वरची अविभाज्य संख्या परत करणारे $p(x)$ हे फलन विचारात घ्या;
म्हणजे $p(0) = 2$, $p(1) = 3$, $p(2) = 5$, इत्यादी. (सोयीसाठी आपण
$p(x)$ हे अनेकदा $p_x$ असे लिहू; $p_0=2$, $p_1=3$, इत्यादी.)

$x$ पेक्षा मोठी असलेली पहिली अविभाज्य संख्या परत करणारे
$\fn{nextPrime(x)}$ हे फलन आपल्याकडे असेल, तर आदिम पुनरावर्तन वापरून
$p$ ची व्याख्या सहज करता येते:
\begin{align*}
  p(0) & = 2\\
  p(x+1) & = \fn{nextPrime}(p(x))
\end{align*}
$\fn{nextPrime}(x)$ ही $y > x$ आणि $y$ अविभाज्य अशी लघुतम संख्या~$y$ असल्यामुळे
तिचे संगणन अपरिबद्ध शोधाने सहज करता येते. पण यूक्लिडच्या एका निकालामुळे
तिची व्याख्या परिबद्ध लघुतमीकरणानेही करता येते: $x$ पेक्षा मोठी आणि
$\fact{x}+1$ पेक्षा मोठी नसलेली अविभाज्य संख्या नेहमीच अस्तित्वात असते.
\[
  \fn{nextPrime(x)} =
  \bmin{y \leq \fact{x}+1}{(y > x \land \fn{Prime}(y))}.
\]
यावरून $\fn{nextPrime}(x)$ आणि म्हणून $p(x)$ ही फलने केवळ संगणनीयच नव्हे,
तर आदिम पुनरावर्तीही आहेत.

(उत्सुकता असेल तर यूक्लिडच्या निकालाचा संक्षिप्त पुरावा असा. दोनपेक्षा
लहान मूल्यांसाठी दोन हीच अपेक्षित अविभाज्य संख्या आहे. इतर मूल्यांसाठी
$p_n$ ही $\le x$ असलेली सर्वांत मोठी अविभाज्य संख्या आहे असे समजा आणि
सर्व $\le x$ अविभाज्य संख्यांचा $p =
p_0\cdot p_1 \cdot \dots \cdot p_n$ हा गुणाकार विचारात घ्या. एकतर $p+1$
अविभाज्य असते, किंवा $x$ आणि~$p+1$ यांच्या दरम्यान एखादी अविभाज्य संख्या
असते. ते का? $p+1$ अविभाज्य नाही असे समजा. मग $q \mid p+1$ अशी एखादी
अविभाज्य संख्या असते, जेथे $q < p+1$. $\le x$ असलेल्या कोणत्याही अविभाज्य
संख्येने $p+1$ ला भाग जात नाही. ($p$ च्या व्याख्येनुसार, प्रत्येक
$p_i \le x$ अविभाज्य संख्येने~$p$ ला भाग जातो, म्हणजे बाकी~$0$ उरते.
म्हणून प्रत्येक $p_i \le x$ अविभाज्य संख्येने $p+1$ ला भागल्यास बाकी~$1$
उरते आणि त्यामुळे $p_i \nmid p+1$.) म्हणून $q$ ही $> x$ आणि $< p+1$
असलेली अविभाज्य संख्या आहे. तसेच $p \le \fact{x}$, म्हणून $> x$ आणि $\le
\fact{x}+1$ असलेली अविभाज्य संख्या अस्तित्वात आहे.)

\begin{prob}
परिबद्ध लघुतमीकरण वापरून पूर्णांक भागाकार $d(x, y)$ ची व्याख्या करा.
\end{prob}

\end{document}
